\problemname{Pianot}
Christian vill lära sig spela piano, men klarar bara av att spela med pekfingret just nu.
Just nu håller han på och över hårt för att lära sig ett stycke.
Ett stycke beskrivs som en sekvens tangenter som ska tryckas i ordning. Exempelvis beskriver
stycket \texttt{abdaa} att han ska trycka på tangenterna
\texttt{a}, \texttt{b}, \texttt{d}, \texttt{a}, \texttt{a} i ordning. 
Du märker dock att han har svårt att röra pekfingret snabbt nog mellan tangenterna.
Märkvärdigt nog tar det honom alltid exakt en sekund att flytta fingret en tanget, men
noll sekunder att spela noten.
För att hjälpa honom tänker du ändra om ordningen på pianots tangenter så att totala tiden
det tar att spela stycket blir så litet som möjligt.

Exempelvis, om stycket han försöker lära sig \texttt{abdaa} och pianot har ordningen
\texttt{abd} \dots kommer han totalt behöva flytta fingret totalt 6 pianotangenters längd.
Avstånden när han spelar stycket blir följande:
\begin{itemize}
  \item \texttt{a} till \texttt{b} tar en sekund.
  \item \texttt{b} till \texttt{d} tar en sekund.
  \item \texttt{d} till \texttt{a} tar två sekunder.
  \item \texttt{a} till \texttt{a} tar noll sekunder.
\end{itemize}
Vilket innebär att det totalt tar fyra sekunder att spela stycket \texttt{abdaa} på ett piano
med ordningen \texttt{abd} \dots.

Skriv ett program som läser ett stycke och skriver ut totala tiden det tar för Christian att
spela stycket på det bästa pianot.

\section*{Indata}
Den första raden innehåller en sträng som består av bokstäverna \texttt{a}-\texttt{z}, proteinet Björn vill göra \texttt{rent}.
Låt antalet bokstäver i proteinet vara $n$. Då gäller alltid $1 \leq n \leq 32$.

På nästa rad följer ett heltal $k$ ($1 \leq k \leq n$), antalet bokstäver Björn kan ta bort med en klippning. 

\section*{Utdata}
Skriv ut ett heltal: det minsta antalet klippningar Björn behöver göra för att proteinet ska bli \texttt{rent}.


\section*{Poängsättning}
Din lösning kommer att testas på en mängd testfallsgrupper.
För att få poäng för en grupp så måste du klara alla testfall i gruppen.

\noindent
\begin{tabular}{| l | l | p{12cm} |}
  \hline
  \textbf{Grupp} & \textbf{Poäng} & \textbf{Gränser} \\ \hline
  $1$    & $20$       & $k = 1$ \\ \hline
  $2$    & $20$       & Proteinet består endast av bokstäverna \texttt{a} och \texttt{b} \\ \hline
  $3$    & $60$       & Inga ytterligare begränsningar. \\ \hline
\end{tabular}
